\documentclass{article}
\usepackage[a4paper, total={6.5in, 10in}]{geometry}
\setlength{\parindent}{0pt}
\usepackage{tcolorbox}
\tcbset{colback=white!94!black, colframe=black!60!white, coltitle = white, boxrule=0.8pt, arc=2mm}
\newtcolorbox{boxs}[2][]{title= #2,#1}
\usepackage{datetime2}
\usepackage[british]{babel}
\usepackage{amsfonts}
\usepackage{amsmath}

\title{\textbf{Linear Algebra HW 3}}
\author{Robert O'Connell}
\date{\today}
\begin{document}
\maketitle
\subsection*{Question 1}
The possibilities for \(i,j\) are 2,6 and for \(k,l\) they are 4,6
\\

An odd permutation is one which has an odd number of inversions. 

We only need to calculate the sign of one of the possible permutations, as then we know if we 
swictch one entry the sign switches.
\\

For the permutation, \(i = 2, \quad  j = 6, \quad k = 4, \quad l = 6\)
\[
\begin{bmatrix}
    3 & 2 & 6 & 1 & 4 & 7 & 5 & 8 \\
    3 & 6 & 8 & 7 & 1 & 5 & 2 & 4
\end{bmatrix}
\]

We then count the number of inversions, and find that there are 7 on top and 17 on the bottom, therefore the inversion is even.
\\

Then if we switch to either just the \(i,j\) entries or just the \(k,l\) entries we will have an odd inversion, as when we swap both, 
two switches of -1 just leave the permutation's sign unchanged.
\\

Thus the values of \(i,j,k,l\) which represent an odd permutation are 
\(2,6,6,4\) and \(6,2,4,6\)

\subsection*{Question 2}
We know that if a matrix is not invertible then the determinant is zero.

Thus we are looking for the values of \(k\) which lead to a zero determinant

\[
A = 
\begin{bmatrix}
    2 & k & k \\
    k & k & k \\
    8 & 7 & k
\end{bmatrix}
\]

Then 
\begin{align*}
    \det{(A)} &= 2 \det{\begin{bmatrix} k & k \\ 7 & k \end{bmatrix}} - k \det{\begin{bmatrix} k & k \\ 8 & k \end{bmatrix}} + k \det{\begin{bmatrix} k & k \\ 8 & 7 \end{bmatrix}} \\
    &= 2(k^2-7k) -k(k^2 -8k) + k(7k-8k) \\
    &= -k^3 + 9k^2 -14k
\end{align*}

\[
k^3 -9k^2 + 14k = 0 
\]
\[
\Rightarrow k = 0, \quad k = 2, \quad k = 7
\]



\subsection*{Question 3}
We need the determinant of
\[
A =
\begin{bmatrix}
0 & 9 & 3 & 9 \\
2 & 2 & 0 & 2 \\
3 & 0 & 1 & 3 \\
5 & -1 & 8 & 7
\end{bmatrix}.
\]

Swap $R_1$ and $R_2$, this changes the determinant of \(A\) by a sign
\[
\begin{bmatrix}
2 & 2 & 0 & 2 \\
0 & 9 & 3 & 9 \\
3 & 0 & 1 & 3 \\
5 & -1 & 8 & 7
\end{bmatrix}.
\]

\(R_1 \times \frac{1}{2}\), this scales the determinant by \(\frac{1}{2}\)
\[
\begin{bmatrix}
1 & 1 & 0 & 1 \\
0 & 9 & 3 & 9 \\
3 & 0 & 1 & 3 \\
5 & -1 & 8 & 7
\end{bmatrix}.
\]

\(
R_3 \to R_3 - 3R_1, \quad
R_4 \to R_4 - 5R_1.
\) This doesn't affect the determinant

\[
\begin{bmatrix}
1 & 1 & 0 & 1 \\
0 & 9 & 3 & 9 \\
0 & -3 & 1 & 0 \\
0 & -6 & 8 & 2
\end{bmatrix}.
\]

\(
R_3 \to R_3 + \tfrac{1}{3} R_2, \quad
R_4 \to R_4 + \tfrac{2}{3} R_2.
\) Similary this doesn't affect the determinant

\[
\begin{bmatrix}
1 & 1 & 0 & 1 \\
0 & 9 & 3 & 9 \\
0 & 0 & 2 & 3 \\
0 & 0 & 10 & 8
\end{bmatrix}.
\]

\(
R_4 \to R_4 - 5R_3.
\)

\[
\begin{bmatrix}
1 & 1 & 0 & 1 \\
0 & 9 & 3 & 9 \\
0 & 0 & 2 & 3 \\
0 & 0 & 0 & -7
\end{bmatrix}.
\]

The determinant of a matrix in upper triangular form is simply the product of the elements along the
main diagonal, but don't forget the changes as a result of the row operations
\[
\det(A) = (1)(9)(2)(-7) \times (2)(-1) = 252.
\]


\subsection*{Question 4}
We need the determinant of
\[
B =
\begin{bmatrix}
2 & 1 & 0 & 3 & -1 \\
0 & 4 & 5 & 1 & 2 \\
1 & 0 & -2 & 2 & 3 \\
3 & -1 & 1 & 0 & 4 \\
2 & 2 & 3 & 1 & 1
\end{bmatrix}.
\]

\(R_1 \times \tfrac{1}{2}\), this scales the determinant by \(\frac{1}{2}\)
\[
\begin{bmatrix}
1 & \tfrac{1}{2} & 0 & \tfrac{3}{2} & -\tfrac{1}{2} \\
0 & 4 & 5 & 1 & 2 \\
1 & 0 & -2 & 2 & 3 \\
3 & -1 & 1 & 0 & 4 \\
2 & 2 & 3 & 1 & 1
\end{bmatrix}.
\]


\(
R_3 \to R_3 - 1\cdot R_1,\quad
R_4 \to R_4 - 3\cdot R_1,\quad
R_5 \to R_5 - 2\cdot R_1,
\) This doesn't effect the determinant

\[
\begin{bmatrix}
1 & \tfrac{1}{2} & 0 & \tfrac{3}{2} & -\tfrac{1}{2} \\[5pt]
0 & 4 & 5 & 1 & 2 \\[5pt]
0 & -\tfrac{1}{2} & -2 & \tfrac{1}{2} & \tfrac{7}{2} \\[5pt]
0 & -\tfrac{5}{2} & 1 & -\tfrac{9}{2} & \tfrac{11}{2} \\[5pt]
0 & 1 & 3 & -2 & 2
\end{bmatrix}.
\]

\(R_2 \times \tfrac{1}{4}\) This scales the determinant by \(\frac{1}{4}\)
\[
\begin{bmatrix}
1 & \tfrac{1}{2} & 0 & \tfrac{3}{2} & -\tfrac{1}{2} \\[5pt]
0 & 1 & \tfrac{5}{4} & \tfrac{1}{4} & \tfrac{1}{2} \\[5pt]
0 & -\tfrac{1}{2} & -2 & \tfrac{1}{2} & \tfrac{7}{2} \\[5pt]
0 & -\tfrac{5}{2} & 1 & -\tfrac{9}{2} & \tfrac{11}{2} \\[5pt]
0 & 1 & 3 & -2 & 2
\end{bmatrix}.
\]

\(
R_3 \to R_3 + \tfrac{1}{2}R_2,\quad
R_4 \to R_4 + \tfrac{5}{2}R_2,\quad
R_5 \to R_5 - 1\cdot R_2,
\)

\[
\begin{bmatrix}
1 & \tfrac{1}{2} & 0 & \tfrac{3}{2} & -\tfrac{1}{2} \\[5pt]
0 & 1 & \tfrac{5}{4} & \tfrac{1}{4} & \tfrac{1}{2} \\[5pt]
0 & 0 & -\tfrac{11}{8} & \tfrac{5}{8} & \tfrac{15}{4} \\[5pt]
0 & 0 & \tfrac{33}{8} & -\tfrac{31}{8} & \tfrac{27}{4} \\[5pt]
0 & 0 & \tfrac{7}{4} & -\tfrac{9}{4} & \tfrac{3}{2}
\end{bmatrix}.
\]

\(
R_4 \to R_4 + 3R_3,\qquad R_5 \to R_5 + \tfrac{14}{11}R_3,
\)

\[
\begin{bmatrix}
1 & \tfrac{1}{2} & 0 & \tfrac{3}{2} & -\tfrac{1}{2} \\[5pt]
0 & 1 & \tfrac{5}{4} & \tfrac{1}{4} & \tfrac{1}{2} \\[5pt]
0 & 0 & -\tfrac{11}{8} & \tfrac{5}{8} & \tfrac{15}{4} \\[5pt]
0 & 0 & 0 & -2 & 18 \\[5pt]
0 & 0 & 0 & -\tfrac{16}{11} & \tfrac{69}{11}
\end{bmatrix}.
\]

\(
R_5 \to R_5 - \tfrac{8}{11}R_4,
\)

\[
\begin{bmatrix}
1 & \tfrac{1}{2} & 0 & \tfrac{3}{2} & -\tfrac{1}{2} \\[5pt]
0 & 1 & \tfrac{5}{4} & \tfrac{1}{4} & \tfrac{1}{2} \\[5pt]
0 & 0 & -\tfrac{11}{8} & \tfrac{5}{8} & \tfrac{15}{4} \\[5pt]
0 & 0 & 0 & -2 & 18 \\[6pt]
0 & 0 & 0 & 0 & -\tfrac{75}{11}
\end{bmatrix}.
\]

The determinant of a matrix in upper triangular form is simply the product of the elements along the
main diagonal, but don't forget the changes as a result of the row operations
\[
\det(B) = \left(1\right)\left(1\right)\left(-\tfrac{11}{8}\right)\left(-2\right)\left(-\tfrac{75}{11}\right) \times (2)(4)
= -150
\]



\subsection*{Question 5}
\[
C^{-1}(A+X)B^{-1} = I_n
\]

Since we know \(C^{-1}C =I_n\), one can see that
\[
(A+X)B^{-1} = C
\]
\[
\Rightarrow A+X = CB
\]
\[
\Rightarrow X = CB - A
\]

\subsection*{Question 6}
This statement is false and can be easily proven with a simple counterexample
\\

Consider the \(2 \times 4 \) augmented matrix of a linear system of two equations in 
three variables

\[
\begin{bmatrix}
    2 & 4 & 6 & 1 \\
    4 & 8 & 12 & 25
\end{bmatrix}
\]

Then when we do the row operation \(r_2 - 2r_1\), we get
\[
\begin{bmatrix}
    2 & 4 & 6 & 1 \\
    0 & 0 & 0 & 23
\end{bmatrix}
\]

There is an obvious inconsistency of \(0 = 23\), which implies the system of linear equations has no solutions

Thus the original statement is false

\subsection*{Question 7}
Since the linear system, with \(\vec{b} = \vec{0}\) has an infinte number of solutions this means that in \(A\) there is not going 
to be a pivot variable in each column, thus there are going to be at least three zero rows when we row reduce \(A\).
\\

Then in the second linear system when we use row reduction and get at least three rows of zero in \(A\), where at least one of them could be equal to a 
non-zero constant, which makes the system inconsistent and have no solutions. 

Or the other possibility is during the row reduction we get an additional zero in \(\vec{b}\) in which
case there are no inconsistencies, still not a pivot variable in each column and thus an infinte number of solutions.

Thus the two possibilities are zero solutions and infinitely many solutions.

\subsection*{Question 8}
We know that the following properites of the determinant of a matrix are true

\begin{itemize}
    \item det\((AB)\) = (det\((A)\))(det\((B)\))
    \item det\((A^T)\) = det\((A)\)
\end{itemize}

Then 
\begin{enumerate}
    \item det\((ABAC)\) = (det\((AB)\)) (det\((AC)\)) = \(\det{(A)}^2\) (det\((B)\)) (det\((C))\) = -200
    \item det\((C^4)\) = det\((CCCC)\) = \(\det{(C)}^4\) = 16
    \item det\((B^TBC)\) = (det\((B^T)\)) (det\((B)\)) (det\((C)\)) = \(\det{(B)}^2\) (det\((C))\) = -32
\end{enumerate}

\subsection*{Question 9}
There are a few facts which are relevant to this question, when you perform the elementray row operation of 
\begin{itemize}
    \item replacement, the determinant of the new matrix remains unchanged
    \item interchanging, the determinant of the new matrix is \(-1\) times the previous determinant
    \item scaling, the determinant of the new matrix is also scaled by the same constant
\end{itemize}


Then the elementray matrices have the following effects on the determinants of the new matrices formed
\begin{itemize}
    \item \(E_1\) has no effect (replacement)
    \item \(E_2\) scales the determinant by \(-\frac{4}{7}\) (scaling)
    \item \(E_3\) changes the sign of the determinant (interchanging)
    \item \(E_4\) has no effect (replacement)
\end{itemize}


Thus
\begin{enumerate}
    \item \(\det{(E_4E_3)} = -1\)
    \item \(\det{(E_4E_1)} = 1\)
    \item \(\det{(E_2E_4)} = -\frac{4}{7}\)
    \item \(\det{(E_1E_2E_3E_4A)} = (1)(-\frac{4}{7})(-1)(1)(5) = \frac{20}{7}\) 
\end{enumerate}

\end{document}